% !PRINTED 465
\weberpage{465}

%% [running head, not body text -- commented out 2026-08-07] Page furniture (running
%% head, folio, signature) is excluded from this edition by standing policy.
%% \noindent \S\,126. \hfill Die Klasseninvarianten $f(\omega)^{24}$. \hfill 465

%<<CONT
1. Wenn $D\equiv 4\pmod{8}$ ist, setzen wir
\[
n=m,\quad x=1,\quad y=0,
\]
\[
\alpha=\frac{B}{2},\quad \beta=A,
\]
\[
\gamma\equiv -C,\quad \delta\equiv\frac{B}{2}\pmod{2}.
\]

Die Kongruenz (9) fordert also, da nicht alle drei Koeffizienten $A$, $B$, $C$ gerade sein können, daß die beiden äußeren Koeffizienten ungerade, der mittlere durch 4 teilbar sei, und dies findet nur für eine der drei Formen (3), (5) statt; folglich genügt eine der drei Funktionen (2) der Gleichung (7), und ist also Klasseninvariante.

\[
\begin{aligned}
2.\quad &D\equiv 0\pmod{8},\\
&n=m+1,\quad x=1,\quad y=\pm 2,\\
&\alpha=\frac{B}{2}-1,\quad \beta=A,\\
&\gamma\equiv -C,\quad \delta\equiv\frac{B}{2}+1\pmod{2}.
\end{aligned}
\]

Es müssen also auch hier die beiden äußeren Koeffizienten ungerade $B\equiv 2\pmod{4}$ sein, und es verhält sich dann alles wie oben. Damit ist bewiesen:

1. \gesperrt{Ist $\omega$ die Wurzel einer primitiven quadratischen Gleichung mit negativer, durch 4 teilbarer Diskriminante $D$ und ungeraden äußeren Koeffizienten, so ist $f(\omega)^{24}$ Klasseninvariante für die Diskriminante $D$.}

Die Annahme, daß $A$ relativ prim zu $n$, d.\,h. zu $m+1$ oder zu $m$ sei, kann nachträglich als unwesentlich aufgegeben werden.

Denn wenden wir auf $\omega$ eine lineare Transformation $(S)$ an, so geht die Gleichung (3) bei ungeraden $A$, $C$ in eine äquivalente Gleichung über, in der die äußeren Koeffizienten dann und nur dann beide ungerade sind, wenn $(S)$ zur ersten oder zweiten Klasse gehört, also wenn $f(\omega)^{24}$ durch $(S)$ ungeändert bleibt. Man kann dann noch $(S)$ so bestimmen, daß der erste Koeffizient in der umgeformten Gleichung (3) zu einer beliebigen Zahl, also auch zu $n$ relativ prim wird.
